\section{Cel, zakres i założenia}
\lead{Cztery współpracujące systemy, jedna sytuacja operacyjna, lokalna zdolność działania podczas awarii.}
Projekt obejmuje monitoring wizyjny CCTV, kontrolę dostępu KD, system sygnalizacji włamania i napadu SSWiN oraz system sygnalizacji pożarowej SSP, określany w temacie także jako SAP. Integracja w centrum nadzoru SOC ma ułatwić ocenę zdarzenia i koordynację reakcji.

Wybrano fikcyjną strefę terminalowo-techniczną lotniska, aby udostępnić opracowanie publicznie i zachować pełną kontrolę nad założeniami. Odniesienie do infrastruktury krytycznej wynika z funkcji transportowej i potrzeby ciągłości usług. Nie przesądza o formalnym zakwalifikowaniu obiektu do IK [1].

\begin{tabularx}{\linewidth}{@{}P{42mm}Y@{}}
\toprule\textbf{Parametr} & \textbf{Założenie koncepcyjne}\\\midrule
Obszar opracowania & 600 $\times$ 400 m = 24 ha; lokalny układ współrzędnych, bez geolokalizacji.\\
Budynki & Terminal 80 $\times$ 40 m; hangar 40 $\times$ 30 m; obiekt techniczny 20 $\times$ 15 m; portiernia 10 $\times$ 8 m.\\
Ruch i obsada & Do 300 osób jednocześnie w terminalu; ochrona obecna całodobowo; dwa stanowiska SOC.\\
Węzły & N0: terminal i SOC; N1: hangar; N2: techniczny; N3: portiernia i dwie bramy.\\
Skala systemu & 24 kamery; 12 drzwi i 2 bramy KD; 24 czujki ruchu i 16 kontaktronów SSWiN; 80 czujek pożarowych.\\
Kontynuacja pracy & LAN z pierścieniem światłowodowym; UPS 30 min; agregat; niezależne baterie central i zasilaczy KD.\\\bottomrule
\end{tabularx}

\subsection*{Granice projektu}
Zakres obejmuje wybrane wejścia, budynki, bramy i rejon płyty postojowej. Nie obejmuje całego obwodu lotniska, pasa startowego, nawigacji, kontroli ruchu lotniczego, ochrony antydronowej ani stref paliwowych. Obsługę bram ograniczono do interfejsu KD; napędy i urządzenia bezpieczeństwa ruchu wymagają odrębnego projektu.

\begin{note}{Co wymaga rozwinięcia przed realizacją}
Brak mapy geodezyjnej, rzeczywistych rzutów, scenariusza pożarowego i kart urządzeń. Lokalizacje, długości i liczby czujek są koncepcją, a nie dokumentacją wykonawczą. Rozwiązania ppoż., zasilanie, ewakuację i wymagania ochrony lotnictwa trzeba uzgodnić w rzeczywistym procesie projektowym.
\end{note}
\src{[1] RCB -- pojęcie infrastruktury krytycznej; pozostałe parametry: założenia własne}
